% !TEX root = thesis-main.tex

\documentclass[12pt,a4paper,openright,twoside]{book}
\usepackage[utf8]{inputenc}
\usepackage{disi-thesis}
\usepackage{code-lstlistings}
\usepackage{notes}
\usepackage{shortcuts}
\usepackage{acronym}

\school{\unibo}
\programme{Corso di Laurea Magistrale in Ingegneria e Scienze Informatiche}
\title{Modernizing a Spark-Based ETL Platform Using Transactional Lakehouse Technologies on AWS}
\author{Tommaso Brini}
\date{\today}
\supervisor{Prof. Boschetti Marco Antonio}
\cosupervisor{Dott. CoSupervisor 1}
\morecosupervisor{Dott. CoSupervisor 2}
\session{II}
\academicyear{2025-2026}
\company{PROMETEIA S.P.A}

% Definition of acronyms
%\acrodef{IoT}{Internet of Thing}
%\acrodef{vm}[VM]{Virtual Machine}


\mainlinespacing{1.241} % line spacing in mainmatter, comment to default (1)

\begin{document}

\frontmatter\frontispiece

\begin{abstract}	
Questa tesi analizza la modernizzazione di una piattaforma ETL basata su Apache Spark e AWS Glue,
precedentemente vincolata a un'architettura Data Lake tradizionale sul formato Parquet.
L'obiettivo centrale del lavoro è l'introduzione del paradigma Transactional Lakehouse tramite l'adozione di Apache Iceberg.
L'introduzione di questo table format ha permesso di abilitare proprietà ACID, funzionalità di time travel e una gestione avanzata dei metadati.

Nel corso del tirocinio è stato progettato e implementato un prototipo in grado di gestire pipeline complesse (INSERT, UPDATE, DELETE e MERGE) 
su dataset di grandi dimensioni, includendo una strategia automatizzata per la sincronizzazione del catalogo e dei metadati.

La soluzione proposta viene quindi confrontata con l'architettura esistente per valutarne gli effetti in termini di prestazioni, 
scalabilità e robustezza operativa, evidenziando vantaggi e limiti delle tecnologie lakehouse nel contesto di pipeline ETL in ambiente cloud.
\end{abstract}

% \begin{dedication} % this is optional
% Optional. Max a few lines.
% \end{dedication}

%----------------------------------------------------------------------------------------
\tableofcontents   
%\listoffigures     % (optional) comment if empty
%\lstlistoflistings % (optional) comment if empty
%----------------------------------------------------------------------------------------

\mainmatter

%----------------------------------------------------------------------------------------
\chapter{Introduzione}
\label{chap:introduction}
%----------------------------------------------------------------------------------------

Il presente lavoro di tesi è stato sviluppato nell'ambito di un tirocinio svolto presso Prometeia S.p.A., azienda attiva nella consulenza e nello sviluppo di soluzioni software per il settore finanziario.

Negli ultimi anni la crescita esponenziale dei dati generati da applicazioni distribuite, sistemi transazionali e servizi cloud ha reso necessaria l'evoluzione delle architetture aziendali per la gestione e l'analisi dei dati.
In questo contesto, i processi di \textbf{Extract, Transform, Load} \textbf{(ETL)} hanno assunto un ruolo centrale nelle pipeline di data engineering, richiedendo soluzioni sempre più scalabili, affidabili e flessibili.

La diffusione dei paradigmi Big Data ha favorito l'adozione dei Data Lake, architetture basate su storage distribuito in grado di memorizzare e gestire grandi quantità di dati eterogenei a costi ridotti. 
Formati colonnari come Apache Parquet hanno quindi rappresentato uno standard de facto per l'elaborazione distribuita tramite framework come Apache Spark. 
Nel modello Lakehouse, i Data Lake moderni si basano infatti su storage cloud economico e formati aperti per supportare workload analitici e di machine learning. \cite{armbrust2021lakehouse} \cite{begoli2021lakehouse}

Nonostante i vantaggi in termini di costo e scalabilità, l'approccio tradizionale basato su file Parquet presenta alcune limitazioni significative.
In particolare, i Data Lake classici non forniscono nativamente proprietà transazionali ACID, una gestione evoluta dei metadati o un supporto efficiente per operazioni di aggiornamento incrementale.
Operazioni come UPDATE, DELETE e MERGE richiedono quindi meccanismi applicativi complessi, spesso basati sulla riscrittura di intere partizioni o dataset, con un aumento della complessità delle pipeline ETL e del rischio di inconsistenze nei dati.

Per superare tali criticità si è diffuso il paradigma del Data Lakehouse, che combina la flessibilità dei Data Lake con caratteristiche tipiche dei Data Warehouse, introducendo transazioni, versionamento dei dati, schema evolution e una gestione avanzata dei metadati. \cite{saha2024disruptor}

Tra le tecnologie emergenti in questo ambito, Apache Iceberg si è affermato come uno dei principali table format open-source per la gestione di dataset analitici su larga scala.
Iceberg introduce un layer di metadata management che separa la struttura logica della tabella dalla disposizione fisica dei file su storage object-based, consentendo la gestione di snapshot consistenti, time travel e operazioni incrementali transazionali. \cite{IcebergDocs}

A differenza di un approccio puramente file-based, Iceberg mantiene lo stato della tabella attraverso metadata versionati e manifest file, consentendo aggiornamenti atomici mediante commit espliciti. \cite{IcebergSpec}

Questa tesi si inserisce nel contesto della modernizzazione di una piattaforma ETL aziendale basata su Apache Spark e servizi AWS.
L'architettura esistente utilizza file Parquet memorizzati su Amazon S3 ed elaborati tramite AWS Glue, presentando criticità legate alla gestione degli aggiornamenti incrementali, alla sincronizzazione del catalogo e alla garanzia di idempotenza dei processi batch.

L'obiettivo del lavoro è progettare e valutare una nuova architettura basata sul paradigma Transactional Lakehouse attraverso l'introduzione di Apache Iceberg.
In particolare, il lavoro analizza l'integrazione del table format con Apache Spark e l'ecosistema AWS, includendo aspetti relativi alla gestione dei metadati, all'esecuzione di operazioni MERGE transazionali e ai processi di maintenance delle tabelle.
Durante il progetto sono inoltre state esplorate le funzionalità offerte da \textbf{Amazon S3 Tables} \cite{S3TablesDocs}, soluzione managed introdotta da AWS per la gestione di tabelle Iceberg.

Il contributo principale della tesi consiste nella progettazione e implementazione di un prototipo in grado di eseguire operazioni incrementali affidabili e idempotenti su dataset di grandi dimensioni, confrontando la nuova soluzione con l'architettura preesistente in termini di prestazioni, scalabilità e robustezza.
L'analisi sperimentale evidenzia inoltre i trade-off introdotti dall'adozione del paradigma lakehouse, valutando vantaggi e limiti sia di un'implementazione Iceberg tradizionale sia dell'approccio gestito offerto da S3 Tables.

Nei prossimi capitoli, la tesi introdurrà lo stato dell'arte relativo alle architetture Data Lake e Lakehouse, descrivendo le principali tecnologie utilizzate (Capitolo 2).
Successivamente, verrà presentata un'analisi dettagliata dell'architettura esistente e delle criticità riscontrate (Capitolo 3), 
seguita dalla progettazione e implementazione della nuova soluzione basata su Apache Iceberg (Capitolo 4).
Infine, i risultati sperimentali e i benchmark effettuati saranno discussi nel Capitolo 5, 
raccogliendo conclusioni e possibili sviluppi futuri nel Capitolo 6.

\chapter{Stato dell'arte e tecnologie}
\label{chap:state-of-the-art}

\begin{itemize}
    \item L'evoluzione delle Architetture Dati: Dal Data Warehouse al Data Lake, fino al paradigma Data Lakehouse.
    \item I Table Format Moderni: Analisi comparativa tra Parquet "puro" e Apache Iceberg.
    \item Focus su Iceberg: Snapshot, Schema Evolution, Partition Evolution e Time Travel.
    \item Apache Spark: Architettura e ruolo nell'elaborazione distribuita.
    \item Ecosistema AWS: Panoramica dei servizi utilizzati (S3, Glue Catalog, AWS Glue Jobs, Athena, S3 Tables).
\end{itemize}

\section{Evoluzione delle Architetture Dati} Medio
Contestualizzare data warehouse, data lake e lakehouse.
Contenuti: (OLTP vs OLAP, separazione compute/storage, nascita dei Data Lake, limiti dei Data Lake tradizionali, nascita del paradigma Lakehouse).
\subsection{Il Data Warehouse Tradizionale}
Carattestiche, vantaggi computazionali (proprietà ACID, indici) e limiti (costi elevati di storage, rigidità computazionale, incapacità di gestire dati non strutturati o semi-strutturati).
\subsection{L'avvento del Data Lake}
La nascita dello storage a basso costo (HDFS, Amazon S3) e dei formati di file aperti (Apache Parquet). 
Il problema del "Data Swamp": mancanza di transazioni, corruzione dei dati in caso di fallimenti, inefficienza negli aggiornamenti incrementali.
\subsection{Paradigma Data Lakehouse}
Come questa nuova architettura riesca a fondere i due mondi, portando l'affidabilità dei data warehouse sopra lo storage economico dei data lake.

\section{Open File Format e Apache Parquet} Medio
Cosa sono gli open file format, perché Parquet è diventato uno standard, quali limiti ha.
Contenuti: row vs column format, compressione, interoperabilità, performance, limiti (assenza ACID, no time travel, update inefficienti).

\section{Apache Spark e l'elaborazione distribuita} Leggero
Obiettivo: contestualizzare il motore ETL utilizzato.
Contenuti: architettura Spark, DataFrame API, Spark SQL, distributed processing, ruolo di Spark nei Data Lake/Lakehouse.

\section{Open Table Format e paradigma Lakehouse}
Obiettivo: introdurre il concetto di table format come evoluzione dei file format, con focus su Apache Iceberg.
Contenuti: differenza (file format vs table format), metadata layer, ACID, CRUD, time travel, schema evolution.

\section{Apache Iceberg: Architettura e funzionalità}
\subsection{Architettura dei metadata}
Descrizione dettagliata del layer di metadata di Iceberg, con focus su Catalog, Metadata File, Manifest List e Manifest File, Snapshot.
\subsection{Operazioni transazionali e Snapshot}
ACID, optimistic concurrency control, MERGE, snapshot isolation e idempotenza.
\subsection{Time Travel e Schema Evolution}
versionamento, query storiche, evoluzione schema, partition evolution.
\subsection{Maintanance e ottimizzazione}
Compaction, snapshot expiration, orphan cleanup, metadata growth.
\subsection{Confronto con Delta Lake e Hudi}
Tabella comparativa tra le tre soluzioni, con focus su metadata management, modalità di scrittura (MOR vs COW), interoperabilità, supporto per partition evolution e compatibilità con i principali motori di elaborazione.

\section{Ecosistema AWS}
\subsubsection{Amazon S3 -> Data Lake Storage}
le caratteristiche di S3 come object storage.
\subsubsection{AWS Glue Catalog}
Il ruolo del catalogo centralizzato per la memorizzazione dei metadati delle tabelle Iceberg
\subsubsection{AWS Glue Jobs}
L'utilizzo di Spark in modalità serverless per l'esecuzione delle pipeline ETL.
\subsubsection{Amazon Athena}
Gli strumenti per l'interrogazione ad-hoc dei dati.
\subsubsection{Amazon S3 Tables}
Nuova soluzione managed di AWS per la gestione di tabelle Iceberg.


\chapter{Analisi As-is: Architettura Esistente}
\label{chap:existing-architecture}

\begin{itemize}
    \item Pipeline ETL esistente: Descrizione del flusso dati basato su Spark e file Parquet.
    \item Pain Points: Analisi dettagliata delle inefficienze riscontrate, i 6 punti critici del documento di formazione.
\end{itemize}

\chapter{Progettazione e Implementazione}
\label{chap:design-and-implementation}

\begin{itemize}
    \item Design della soluzione (To-Be): Come cambia l'architettura logica.
    \item Configurazione dello Stack: Integrazione di Iceberg con Spark su AWS.
    \item Ottimizzazione: Implementazione di processi di Maintenance (Compaction dei file, rimozione dei vecchi snapshot/orphan files).
\end{itemize}

\chapter{Test e risultati}
\label{chap:testing-and-results}

\begin{itemize}
    \item Benchmark di Performance: Confronto tra il vecchio sistema e il nuovo (tempi di esecuzione del batch).
    \item Affidabilità: Dimostrazione del supporto alle transazioni ACID (es. cosa succede se un job fallisce a metà).
    \item Costi: Impatto dell'ottimizzazione dei file (Compaction) sui costi di storage.
    \item Risposta ai 6 punti critici
\end{itemize}


\chapter{Conclusioni e sviluppi futuri}

\begin{itemize}
    \item Risultati raggiunti: Sintesi dei benefici ottenuti dall'azienda.
    \item Limiti riscontrati: Eventuali criticità del framework o dei servizi AWS durante lo sviluppo.
    \item Evoluzioni: aggiornamento parziale sui record compositi
\end{itemize}
%You may also put some code snippet (which is NOT float by default), eg: \cref{lst:random-code}.

%\lstinputlisting[float,language=Java,label={lst:random-code}]{listings/HelloWorld.java}

%\section{Fancy formulas here}

%----------------------------------------------------------------------------------------
% BIBLIOGRAPHY
%----------------------------------------------------------------------------------------

\backmatter

\nocite{*} % Remove this as soon as you have the first citation

\bibliographystyle{alpha}
\bibliography{bibliography}

\begin{acknowledgements} % this is optional
Optional. Max 1 page.
\end{acknowledgements}

\end{document}
